\section{The pair of shifts \texorpdfstring{$(\frac13,\frac23)$}{(1/3,2/3)}}
\label{sec:chi3-quarter}

For the function $F_\chi$ of \eqref{eq:frames} we use the alternating geometry of parameter $c=\frac12$ at all even scales
$r$. Then $\alpha=\frac14$ and $m=\lceil r/4-\frac12\rceil$. The decay and the nonvanishing are given by the general
Theorems~\ref{thm:decay} and~\ref{thm:selection}. What changes is the content. The function has three digamma terms. So at a
prime $p$ the increments $g(k)$ of \S\ref{sec:construction} have poles in three residue classes. The local bounds of
Section~\ref{sec:content} then become expressions in two first poles instead of one. We first describe these poles. We
then give the bounds that we use for $F_\chi$, and finally the constants.

\subsection{The first poles}
Let $p>3$ be a prime. The increment $g(k)=\frac1{3k-2}-\frac1{2k-1}+\frac1{6k}$ has a pole modulo $p$ exactly when $p$
divides $k$, $3k-2$ or $2k-1$. Let $k_A$ and $k_B$ be the least $k\ge1$ with $p\mid3k-2$ and with $p\mid2k-1$. Then
\[
 k_B=\frac{p+1}2,\qquad
 k_A=\begin{cases}\frac{p+2}3,&p\equiv1\pmod6,\\[3pt]\frac{2p+2}3,&p\equiv5\pmod6.\end{cases}
\]
The residues of $p\,g(k)$ at $k_A$ and $k_B$ are $1$ and $-1$ if $p\equiv1\pmod6$, and $\frac12$ and $-1$ if
$p\equiv5\pmod6$. In the first case $k_A<k_B$ and the two residues cancel. So the partial sums $\sum_{k\le t}g(k)$ are
$p$-integral for $t<k_A$, and again for $k_B\le t<p$.

Both first poles are affine in one index. Define $k_p$ by $6k_p=jp+5$ with $0<j<6$. Then $j=1$ for $p\equiv1\pmod6$ and
$j=5$ for $p\equiv5\pmod6$, $k_p=\frac j6p+O(1)$, and
\[
 (k_A,k_B)=\begin{cases}(2k_p-1,\ 3k_p-2),&p\equiv1\pmod6,\\[2pt](2k_p-p-1,\ 3k_p-2p-2),&p\equiv5\pmod6.\end{cases}
\]
Thus $j$ plays the role of the root class of \S\ref{sec:content}, with $q$ replaced by $6$.

\subsection{The local bounds}
Fix a prime $p>6$ with $6\bar N<p^2$. Every pole below $\bar N$ is then simple. Call a node $t$ \emph{regular} if
$t<\min(k_A,k_B)$, or if $p\equiv1\pmod6$ and $t$ is a border node with $k_B\le t<p$. At a regular node the partial sum
of the increments up to $t$ is $p$-integral. When $t$ is a common node, its diagonal entry is $p$-integral too.

\begin{proposition}\label{prop:chi3-local}
Every member of the alternating geometry of parameter $\frac12$ for $F_\chi$ satisfies $e_p(r,I)\le\tau$ for each of the
following $\tau$.
\begin{enumerate}[label=(\roman*),leftmargin=*]
\item The bounds (T1) and (T1$'$), namely $2(r+m)$ and $2r+2m-\|m\|_p$.
\item If $\bar N\le2p$, the bound $2r+2m-n_{\rm reg}$. Here $n_{\rm reg}$ is the number of pairs formed by a regular node $t$
 and a side on which $t$ lies, row or column, such that $t+p$ is not a node of the other side.
\item If $\bar N<\min(k_A,k_B)$, the bound $0$.
\item The bound $2r+2m-(t_2-t_1)$ of Theorem~\ref{thm:window}, for every $n$ and every window $[t_1,t_2)$ as there, with
 the cuts $t\equiv0$, $t\equiv k_A$, $t\equiv k_B$ and $t\equiv r\pmod p$.
\end{enumerate}
\end{proposition}

\begin{proof}
The bounds (i) are proved as in \S\ref{sssec:trop}. For (ii) we apply Lemma~\ref{lem:trop} with the following potentials.
At a node that is not regular, the potential is $1$ on its row and on its column. At a regular node $t$, the potential on
its row is $1$ if $t+p$ is a column node, and $0$ otherwise. The potential on its column is $1$ if $t+p$ is a row node,
and $0$ otherwise. Their sum is $2(r+m)-n_{\rm reg}$, and they dominate the weights $W$. Indeed, a diagonal weight is $0$
at a regular common node and at most $2$ in general. Two nodes that are not congruent modulo $p$ have weight at most $1$,
since all poles are simple. Their weight is $0$ if both are regular. Two distinct congruent nodes have weight at most
$2$. Since $\bar N\le2p$ and regular nodes are smaller than $p$, at most one of them is regular. Suppose that the row node $t$
of such an entry is regular. Then its column node is $t+p$, so the potential on the row of $t$ is $1$. The case of a
regular column node is symmetric.

For (iii), the condition $\bar N<\min(k_A,k_B)$ implies $\bar N<p$. Then no index below $\bar N$ is a pole and no two nodes are
congruent modulo $p$, so every entry is $p$-integral. For (iv), the proof of Theorem~\ref{thm:window} uses $F$ only
through its poles. The step function $V$ and the diagonal residues are constant between consecutive cuts. This is what
the translation between the fibres of the window needs.
\end{proof}

The bounds (T2)--(T5), (LP1), (LP2), (R) and (S) of Theorem~\ref{thm:local} are not used for $F_\chi$.

\subsection{The content}
The checker of \S\ref{ssec:checker} and Theorem~\ref{thm:content} apply to $F_\chi$ with $q$ replaced by $6$. The root
class is $j\in\{1,5\}$ as above, and the bounds are those of Proposition~\ref{prop:chi3-local}. These bounds are
expressions of the class $\mathcal E$ in $r,p,m,k_A,k_B$. On each class of primes modulo $6$ the substitution above makes
them expressions in $r,p,m,k_p$. Their homogenisations are taken with $\kappa=j/6$. The proof of
Theorem~\ref{thm:content} uses $F$ only through these bounds and the root class. It goes through with the bound
(iii) in place of (T5) beyond the end of the table.

The table of the formal proof starts at $\theta_0=\frac1{68}=\alpha/17$. Below it the trivial bound $\ell_0r$ with
$\ell_0=\frac52=2+2\alpha$ is used. On $(\alpha/17,\alpha]$ the table uses the bounds (i) and the windows (iv) with
$1\le n\le16$. Beyond $\alpha$ it uses (i), (ii), the windows (iv) with $n=0$ between consecutive cuts, and the bound
(iii), which ends the table. The checker accepts the table for the two classes $j=1,5$. The arithmetic estimate therefore
holds for every even $r$ and every member with $|I|\le64$, with the denominator exponent
\[
 \Lambda_\chi=4.09592776731\ldots
\]
This is an explicit rational number, whose numerator and denominator have $117$ and $116$ digits.

\subsection{Decay and nonvanishing}
The decay estimate is Theorem~\ref{thm:decay} for the function $F_\chi$, with $|\mathcal T|=3$ and $c=\frac12$. The node
measure $\nu_{1/2}$ has mass $\frac54$ and energy
\[
 \cI(\nu_{1/2})=\frac{75}{32}-\frac98\log\frac32-\frac1{16}\log2 .
\]
A certificate measure on a polygon, with a cover by cells accepted by the checks of \S\ref{ssec:slope-cert}, gives
\[
\begin{gathered}
 D=\frac{4643511791526209475773}{10^{21}}=4.64351\ldots,\\ \gamma=\frac{37521870783566737}{5\cdot10^{15}}=7.50437\ldots
\end{gathered}
\]
Nonvanishing is Theorem~\ref{thm:selection} for $F_\chi$. Its increment $\Delta_F$ has its pole at $-1$ isolated in its
class modulo $\Z$, and the depth is $K\le64$. The deduction of Section~\ref{sec:proof}, at the even scales, gives
\[
 \mu\bigl(\Psi(\tfrac13,\tfrac23)\bigr)\le\frac\gamma{D-\Lambda_\chi}=\MuChiExact\ldots<\MuChi,
\]
since $\Psi(\frac13,\frac23)=36\,\xi_\chi$.
